%% ============================================================================
%%  Simulation Campaign Plan
%%  Fragmentation of Compressed Disks as a Phase Transition (VCEM)
%%  N. M. Ghoniem
%% ============================================================================
\documentclass[11pt]{article}
\usepackage[margin=1in]{geometry}
\usepackage{amsmath,amssymb,bm}
\usepackage{booktabs,longtable,array}
\usepackage{xcolor}
\usepackage[colorlinks=true,linkcolor=blue,citecolor=blue,urlcolor=blue]{hyperref}
\usepackage{enumitem}
\setlist{nosep,leftmargin=1.5em}

\newcommand{\Pinf}{P_\infty}
\newcommand{\xc}{x_c}
\newcommand{\rhoc}{\rho_c}
\newcommand{\rhostar}{\rho^{\star}}

\title{\textbf{A Tractable, Self-Learning Simulation Campaign for\\
Fragmentation as a Phase Transition with VCEM}}
\author{N. M. Ghoniem}
\date{\today}

\begin{document}
\maketitle

\begin{abstract}
This document specifies a complete, budget-bounded plan for the VCEM simulation
campaign that will populate the results section of the manuscript
\emph{``Fragmentation of Compressed Disks as a Phase Transition.''} It articulates the
logic behind every parametric choice, fixes the ranges and combinations of input
parameters, sizes the ensembles, and reconciles the total cost against a single
workstation running continuously for 10--20 days. It then provides a master table of
runs that can be executed with agentic assistance, and defines a self-learning
controller that reads the results of completed runs, measures convergence of each
target quantity, and reallocates the remaining compute toward the runs that most
reduce the uncertainty in the paper's headline results.
\end{abstract}

\tableofcontents

%% ============================================================================
\section{Objectives and the statistics the paper requires}
\label{sec:objectives}
%% ============================================================================

The campaign exists to measure a small number of well-defined quantities. Every run
is justified by its contribution to one of them. Table~\ref{tab:targets} maps each
paper deliverable to the observable that produces it and to the design feature that
the observable demands. The guiding principle is that \emph{the control parameter is
internal to each run}: as the load increases, a single VCEM simulation traverses the
entire path from an intact disk to a fragmented one, so one run yields the whole
curves $\Pinf(x)$, $\chi(x)$, $\xi(x)$, the spanning/peak/fragmentation thresholds,
and a complete avalanche sequence. Consequently we do \emph{not} spend separate runs
per value of the control parameter; we spend runs on \emph{ensembles} (random seeds)
and on \emph{system sizes}, and we harvest the rest from each trajectory.

\begin{table}[h]
\centering
\small
\renewcommand{\arraystretch}{1.25}
\begin{tabular}{p{4.6cm} p{4.8cm} p{4.6cm}}
\toprule
\textbf{Paper deliverable} & \textbf{Observable / estimator} & \textbf{Design requirement}\\
\midrule
Critical point(s) $\xc$ ($x_{\rm span}$, $x_{\rm peak}$, $x_{\rm frag}$)
 & peak of $\chi_S(x)$; threshold crossings of $\Pinf$, $R(x)$, $F$
 & ensemble at each size; fine load increments near the peak \\
Static exponents $\beta,\gamma,\nu$
 & finite-size scaling of $\Pinf$, $\chi_{S,\max}$, $\xi$; data collapse
 & $\ge 4$ system sizes spanning a factor $\sim$4 in $L$ \\
Cluster exponent $\tau$
 & fit of $n_s\sim s^{-\tau}$ at $\xc$
 & large-size runs at criticality (tail statistics) \\
Avalanche exponent $\tau_A$
 & MLE fit of $P(A)\sim A^{-\tau_A}$ (energy drops)
 & pooled over \emph{all} runs; essentially free \\
Order of the transition (continuous / first-order / hybrid)
 & shape of $\Pinf$ distribution; hysteresis; $\chi$ divergence vs.\ jump
 & disorder-strength sweep; replicate ensembles \\
Universality verdict
 & comparison of $\{\beta,\gamma,\nu,\tau,\tau_A\}$ to reference classes
 & all of the above, with bootstrap confidence intervals \\
Protocol \& anisotropy robustness
 & shift of $\xc$ and exponents under load/displacement, orientation
 & small control ensembles \\
\bottomrule
\end{tabular}
\caption{Mapping from paper deliverables to observables and to the design features
that the campaign must provide.}
\label{tab:targets}
\end{table}

%% ============================================================================
\section{Logic of the parametric selection}
\label{sec:logic}
%% ============================================================================

\subsection{Which parameters matter, and why}

\paragraph{System size $L$ (primary, expensive).}
Finite-size scaling is the only reliable route to the critical exponents and to a
size-independent $\xc$. It requires several sizes spanning as wide a range as the
budget allows. The cost penalty is severe: at fixed crack density the initial flaw
count grows as $N_0 \propto L^2$, and the dense VCEM solve costs roughly
$\mathcal{O}(N_0^2)$ per increment with $\mathcal{O}(N_0)$ increments, so run time
scales steeply with $L$. Size is therefore the axis on which we spend most carefully,
using a \emph{pyramid} of ensembles (many cheap small runs, few decisive large runs).

\paragraph{Initial crack density $\rhoc$ (control / phase-boundary axis).}
Density sets the proximity of the initial state to the connectivity threshold and is
the natural ``occupation probability'' analog. Holding $\rhoc$ fixed isolates
finite-size effects; sweeping $\rhoc$ at fixed $L$ maps the phase boundary and tests
whether damage, load, and density are equivalent control parameters.

\paragraph{Disorder strength $w$ (transition-order axis).}
The width of the flaw-length (or threshold) distribution is the parameter known from
statistical-fracture theory to control whether breakdown is continuous and
percolation-like (strong disorder) or abrupt and first-order (weak disorder).
Resolving Hypotheses I--IV of the manuscript requires at least three disorder levels.

\paragraph{Flaw orientation (anisotropy axis).}
Aligned flaw populations bias coalescence and can shift $\xc$ and, possibly, the
exponents. A three-level isotropic$\rightarrow$aligned sweep bounds this effect.

\paragraph{Loading protocol (robustness).}
Load- vs.\ displacement-control must give the same intensive critical behavior if the
transition is genuine; a single paired ensemble checks this.

\paragraph{Fixed quantities.}
Elastic constants $E,\nu$ and toughness $G_c$ set only scales and are held fixed.
Crack-growth rule (maximum energy-release-rate, mixed-mode) and the connectivity
tolerance are fixed across the campaign so that all runs are comparable.

\subsection{Ranges and combinations}

\begin{itemize}
\item \textbf{Sizes:} four nominal sizes with initial flaw counts
$N_0 \in \{64,\,256,\,576,\,1024\}$, i.e.\ linear sizes in the ratio
$1\!:\!2\!:\!3\!:\!4$ (since $N_0\propto L^2$). An optional fifth anchor
($N_0\approx 2304$, ratio 6) is held in reserve for the adaptive loop.
\item \textbf{Density:} $\rhoc/\rhostar \in \{0.50,0.70,0.85,1.00,1.15,1.30,1.50\}$,
where $\rhostar$ is the reference density chosen (from pilots) near the expected
critical density.
\item \textbf{Disorder:} length-distribution coefficient of variation
$w \in \{0.15\ (\text{narrow}),\,0.40\ (\text{medium}),\,0.80\ (\text{broad})\}$.
\item \textbf{Orientation:} $\{$isotropic, $30^\circ$ half-width aligned, strongly
aligned$\}$.
\item \textbf{Protocol:} $\{$displacement-control (default), load-control$\}$.
\item \textbf{Dominant-flaw checks:} dominant-to-background length ratio
$r \in \{2,4\}$ on a medium background.
\item \textbf{Seeds:} random realizations per cell, allocated as a size-pyramid in the
finite-size campaign and flat ($\sim$30--40) elsewhere.
\end{itemize}

A full factorial over these axes is neither necessary nor affordable. Instead we use a
\emph{one-reference, several one-axis sweeps} design: a single reference cell
($\rhostar$, medium disorder, isotropic, displacement-control) is shared by every
campaign, and each campaign perturbs exactly one axis away from it. This is the
classical efficient screening design; it isolates each effect, reuses the reference
ensemble, and keeps the run count near 700 rather than the many thousands a factorial
would demand.

%% ============================================================================
\section{Cost model and budget reconciliation}
\label{sec:budget}
%% ============================================================================

\subsection{Per-run cost model}

Calibrate the wall-clock time per run against the initial flaw count with a two-parameter
power law,
\begin{equation}
t(N_0) \;=\; t_{\rm ref}\left(\frac{N_0}{N_{\rm ref}}\right)^{p},
\label{eq:cost}
\end{equation}
where $p\approx 1.5$--$2.0$ reflects the dense-solve scaling. The planning estimates
below assume the representative times in Table~\ref{tab:times}; \textbf{these must be
recalibrated from 3--5 pilot runs} (Section~\ref{sec:pilot}) and Eq.~\eqref{eq:cost}
re-evaluated before committing the budget.

\begin{table}[h]
\centering
\small
\renewcommand{\arraystretch}{1.2}
\begin{tabular}{lccc}
\toprule
Size label & $N_0$ (flaws) & linear ratio & planning time/run \\
\midrule
S  & 64   & 1 & 0.20 h ($\sim$12 min) \\
M  & 256  & 2 & 0.80 h ($\sim$48 min) \\
L  & 576  & 3 & 2.5 h \\
XL & 1024 & 4 & 6.0 h \\
(XXL, reserve) & 2304 & 6 & $\sim$20 h \\
\bottomrule
\end{tabular}
\caption{Representative per-run times used for planning. Recalibrate with pilots.}
\label{tab:times}
\end{table}

\subsection{Total budget and calendar mapping}

The baseline design (Section~\ref{sec:table}) totals $\approx 706$ runs and
$\approx 620$ run-hours of serial-equivalent compute. We add a $\sim$20\,\%
\emph{adaptive reserve} ($\approx 130$ run-hours) that the self-learning controller
spends on refinement, bringing the plan to $\approx 750$ run-hours. The calendar time
depends on how many runs execute concurrently on the workstation
(Table~\ref{tab:calendar}). Because small and medium runs are the bulk of the count,
the recommended mode is to run several of them in parallel (pinned to disjoint core
sets) while reserving more threads for the few L/XL runs.

\begin{table}[h]
\centering
\small
\renewcommand{\arraystretch}{1.2}
\begin{tabular}{cccc}
\toprule
Concurrent runs & Throughput (run-h/day) & Baseline 620 h & With reserve 750 h \\
\midrule
1 (serial)   & 24  & 25.8 d & 31.3 d \\
2            & 48  & 12.9 d & 15.6 d \\
3            & 72  & 8.6 d  & 10.4 d \\
4            & 96  & 6.5 d  & 7.8 d  \\
6            & 144 & 4.3 d  & 5.2 d  \\
\bottomrule
\end{tabular}
\caption{Calendar time vs.\ concurrency. \textbf{Recommendation:} 2--3 concurrent runs
fits the 10--20 day window with the full adaptive reserve; 4--6 finishes the baseline
early and frees days for larger-$L$ anchors and critical-region refinement.}
\label{tab:calendar}
\end{table}

%% ============================================================================
\section{Ensemble allocation strategy}
\label{sec:ensembles}
%% ============================================================================

Statistical precision on a finite-size-scaling fit is dominated by the cheap small
sizes (which can be run many times) anchored by a few expensive large sizes (which fix
the asymptotic slope). We therefore allocate seeds as a \emph{pyramid} in the
finite-size campaign:
\[
\text{S}:160,\quad \text{M}:80,\quad \text{L}:40,\quad \text{XL}:16 .
\]
All other campaigns use flat ensembles of 30--40 seeds, which suffice because they ask
only for the \emph{shift} of an already-located transition, not for an absolute
exponent. The reference cell (size M, $\rhostar$, medium disorder, isotropic,
displacement) is computed once and \emph{reused} by every campaign that includes it,
eliminating redundant runs.

%% ============================================================================
\section{Master run table}
\label{sec:table}
%% ============================================================================

Table~\ref{tab:runs} is the executable plan. Each row is a \emph{cell}: a fixed
parameter combination replicated over the listed number of random seeds. The agent
executes cells in priority order (P0 first), and within a cell executes seeds until the
listed count is reached or the controller (Section~\ref{sec:learning}) declares the
relevant target converged. ``Reuse'' marks cells already produced by Campaign~1.

\begin{longtable}{l l c c c c c c c c}
\caption{Master run table. $N_0$ = initial flaws; $w$ = disorder (CoV); seeds =
realizations; h/run from Table~\ref{tab:times}; cell-h = seeds$\times$h/run.}
\label{tab:runs}\\
\toprule
ID & Campaign & $N_0$ & $\rhoc/\rhostar$ & $w$ & Orient. & Protocol & Seeds & h/run & Cell-h \\
\midrule
\endfirsthead
\multicolumn{10}{c}{\tablename~\ref{tab:runs} (continued)}\\
\toprule
ID & Campaign & $N_0$ & $\rhoc/\rhostar$ & $w$ & Orient. & Protocol & Seeds & h/run & Cell-h \\
\midrule
\endhead
\midrule
\multicolumn{10}{r}{\emph{continued on next page}}\\
\endfoot
\bottomrule
\endlastfoot
\multicolumn{10}{l}{\textbf{Campaign 1 --- Finite-size scaling (P0, primary)}}\\
C1-S  & FSS & 64   & 1.00 & 0.40 & iso & disp & 160 & 0.20 & 32.0 \\
C1-M$^\dagger$ & FSS & 256  & 1.00 & 0.40 & iso & disp & 80  & 0.80 & 64.0 \\
C1-L  & FSS & 576  & 1.00 & 0.40 & iso & disp & 40  & 2.50 & 100.0 \\
C1-XL & FSS & 1024 & 1.00 & 0.40 & iso & disp & 16  & 6.00 & 96.0 \\[2pt]
\multicolumn{10}{l}{\textbf{Campaign 2 --- Density sweep (P1)}}\\
C2-a & dens & 256 & 0.50 & 0.40 & iso & disp & 30 & 0.80 & 24.0 \\
C2-b & dens & 256 & 0.70 & 0.40 & iso & disp & 30 & 0.80 & 24.0 \\
C2-c & dens & 256 & 0.85 & 0.40 & iso & disp & 30 & 0.80 & 24.0 \\
C2-d & dens & 256 & 1.00 & 0.40 & iso & disp & \multicolumn{3}{c}{reuse C1-M} \\
C2-e & dens & 256 & 1.15 & 0.40 & iso & disp & 30 & 0.80 & 24.0 \\
C2-f & dens & 256 & 1.30 & 0.40 & iso & disp & 30 & 0.80 & 24.0 \\
C2-g & dens & 256 & 1.50 & 0.40 & iso & disp & 30 & 0.80 & 24.0 \\[2pt]
\multicolumn{10}{l}{\textbf{Campaign 3 --- Disorder strength (P1)}}\\
C3-narrow & dis & 256 & 1.00 & 0.15 & iso & disp & 40 & 0.80 & 32.0 \\
C3-med    & dis & 256 & 1.00 & 0.40 & iso & disp & \multicolumn{3}{c}{reuse C1-M} \\
C3-broad  & dis & 256 & 1.00 & 0.80 & iso & disp & 40 & 0.80 & 32.0 \\[2pt]
\multicolumn{10}{l}{\textbf{Campaign 4 --- Orientation / anisotropy (P2)}}\\
C4-iso     & ori & 256 & 1.00 & 0.40 & iso     & disp & \multicolumn{3}{c}{reuse C1-M} \\
C4-partial & ori & 256 & 1.00 & 0.40 & $30^\circ$ & disp & 30 & 0.80 & 24.0 \\
C4-aligned & ori & 256 & 1.00 & 0.40 & aligned & disp & 30 & 0.80 & 24.0 \\[2pt]
\multicolumn{10}{l}{\textbf{Campaign 5 --- Protocol \& dominant-flaw checks (P2)}}\\
C5-load    & prot & 256 & 1.00 & 0.40 & iso & load & 30 & 0.80 & 24.0 \\
C5-domin2  & dom  & 256 & 1.00 & 0.40 & iso & disp & 30 & 0.80 & 24.0 \\
C5-domin4  & dom  & 256 & 1.00 & 0.40 & iso & disp & 30 & 0.80 & 24.0 \\
\midrule
\multicolumn{9}{l}{\textbf{Baseline total (new runs only)}} & \textbf{$\approx$620 h}\\
\multicolumn{9}{l}{Adaptive reserve ($\sim$20\,\%)} & $\approx$130 h\\
\multicolumn{9}{l}{\textbf{Planned grand total}} & \textbf{$\approx$750 h}\\
\end{longtable}

\noindent$^\dagger$ C1-M is the shared \emph{reference cell}; campaigns 2--5 reuse it
for their isotropic / medium-disorder / displacement / $\rhostar$ point.

Baseline new-run counts by campaign: C1 $=296$, C2 $=180$, C3 $=80$, C4 $=60$,
C5 $=90$; total $\approx 706$ runs.

%% ============================================================================
\section{Self-learning adaptive controller}
\label{sec:learning}
%% ============================================================================

The campaign is not a static list executed blindly; it is a feedback loop. After every
completed batch (e.g.\ daily, or every $k$ runs) an aggregator recomputes the target
quantities and their uncertainties, and a controller decides where the next runs go.
This converts a fixed budget into the \emph{most informative} set of runs achievable
within it.

\subsection{Per-run record}

Each run writes one machine-readable record (JSON). The controller never re-parses raw
fields; it reads these summaries.
\begin{verbatim}
{
  "run_id": "C1-L_seed037",
  "campaign": "FSS",
  "params": {"N0":576,"rho_ratio":1.00,"w":0.40,
             "orient":"iso","protocol":"disp","seed":37,
             "dominant_ratio":null},
  "status": "completed",          "wall_time_h": 2.43,
  "x_tip":0.31,"x_span":0.62,"x_peak":0.66,"x_frag":0.71,
  "chi_max":184.2,"x_at_chi_max":0.645,
  "Pinf_curve":"arrays/C1-L_seed037_Pinf.npy",
  "xi_curve":"arrays/C1-L_seed037_xi.npy",
  "ns_at_xc":"arrays/C1-L_seed037_ns.npy",
  "avalanches":"arrays/C1-L_seed037_dE.npy",
  "n_avalanches":4172,"tauA_run":1.9,
  "N_fragments":7,"Pinf_final":0.88,
  "converged_flags":{"solver":true,"increment_refined":true}
}
\end{verbatim}

\subsection{Aggregated metrics (recomputed each batch)}

\begin{itemize}
\item Per size $L$: ensemble-mean $\langle\Pinf(x)\rangle$, $\langle\chi_S(x)\rangle$,
$\langle\xi(x)\rangle$; the peak height $\chi_{S,\max}(L)$ and its location
$x_c(L)$, each with a \emph{bootstrap} confidence interval over seeds.
\item Finite-size-scaling fits for $\beta/\nu$, $\gamma/\nu$, $\nu$, and the
data-collapse quality $Q$ (residual of the collapse onto the master curve).
\item Pooled cluster fit $\tau$ at $\xc$ and pooled avalanche fit $\tau_A$ by maximum
likelihood with a Clauset--Shalizi--Newman lower-cutoff selection and a
Kolmogorov--Smirnov goodness-of-fit $p$-value.
\item Transition-order diagnostics: the across-seed distribution of $\Pinf$ near $\xc$
(unimodal vs.\ bimodal), and any hysteresis between loading and unloading.
\end{itemize}

\subsection{Decision rules}

After each batch the controller applies the following rules, in order.
\begin{enumerate}
\item \textbf{Convergence test.} For each target quantity compute the relative
confidence-interval half-width. Targets: $x_c$ to $\le 5\%$, exponent ratios to
$\le 10\%$, $\tau_A$ to $\le 10\%$. Mark converged targets and freeze the cells that
only serve them.
\item \textbf{Greedy variance reduction.} Direct the next seeds to the \emph{unfrozen}
cell with the largest expected reduction in a target's variance \emph{per CPU-hour}
(a budget-aware bandit rule). Because variance falls as $1/n_{\rm seed}$ and cost is
fixed per cell, this favors cheap cells until their marginal value drops below that of
an expensive cell.
\item \textbf{Critical-region refinement.} If the sampled $\chi_S(x)$ peak is resolved
by fewer than $\sim$8 load increments across its full width at half maximum, switch new
runs in that cell to adaptive load stepping (finer $\Delta x$ near $x_{\rm peak}$)
rather than adding seeds.
\item \textbf{FSS adequacy.} Fit the exponents; if the collapse quality $Q$ is poor or
the exponents drift monotonically with $L$ (uncorrected scaling corrections),
\emph{spend the reserve} on the XXL anchor ($N_0\approx2304$) to extend the lever arm.
\item \textbf{Transition-order branch.} If Campaign~3 shows a bimodal $\Pinf$ or
hysteresis at any disorder level, spawn a focused sub-ensemble (extra seeds + finer
increments) at that level to characterize the jump/spinodal (Hypotheses III--IV).
\item \textbf{Budget governor.} Track consumed vs.\ remaining run-hours and projected
calendar. Keep $\ge$15\,\% reserve until P0 and P1 are converged. If behind schedule,
shed work from the lowest priority upward (C5, then C4), never from C1.
\end{enumerate}

\subsection{Stopping criteria}

Stop when either (a) all P0 and P1 targets are converged to tolerance and the P2 checks
have at least their baseline seeds, or (b) the budget (with reserve) is exhausted. On
stop, emit the final aggregated tables and figures (Section~\ref{sec:analysis}).

%% ============================================================================
\section{Execution protocol with agentic help}
\label{sec:execution}
%% ============================================================================

\paragraph{Directory layout.}
\texttt{runs/<run\_id>/} holds \texttt{input.json}, \texttt{output.json},
\texttt{solver.log}, and \texttt{arrays/}. A top-level \texttt{manifest.csv} lists every
cell, its target seed count, completed seeds, and status. A \texttt{controller\_state.json}
records frozen cells, remaining budget, and the current decision.

\paragraph{Agent loop.}
(1) Read \texttt{manifest.csv} and \texttt{controller\_state.json};
(2) select the next \texttt{(cell, seed)} per the priority and variance rules;
(3) write \texttt{input.json} and launch the VCEM solver;
(4) on completion, parse outputs into \texttt{output.json} and append to the aggregates;
(5) re-run the controller; (6) write the daily report; (7) repeat. The loop is
checkpointed after every run so it resumes cleanly after an interruption or reboot.

\paragraph{Pilot calibration.}
\label{sec:pilot}
Before the main campaign, run one seed at each of S, M, L (3--5 runs total), measure
$t(N_0)$, fit Eq.~\eqref{eq:cost}, and rescale Tables~\ref{tab:times}--\ref{tab:calendar}.
Also confirm that $\rhostar$ brackets the transition (i.e.\ that the reference cell
actually reaches spanning and fragmentation); if not, shift $\rhostar$ and the
density grid before launching.

%% ============================================================================
\section{Analysis pipeline and mapping to the paper}
\label{sec:analysis}
%% ============================================================================

\begin{itemize}
\item \textbf{$\xc$ and exponent ratios:} locate $x_c(L)$ from $\chi_{S,\max}(L)$;
fit $\chi_{S,\max}\sim L^{\gamma/\nu}$, $\Pinf(L,\xc)\sim L^{-\beta/\nu}$, and the
collapse $\Pinf L^{\beta/\nu}=\mathcal{F}[(x-\xc)L^{1/\nu}]$ to obtain
$\beta/\nu,\gamma/\nu,\nu$ with bootstrap CIs. (Manuscript Eqs.\ for finite-size
scaling and data collapse.)
\item \textbf{$\tau$:} fit $n_s\sim s^{-\tau}$ pooled over the largest sizes at $\xc$.
\item \textbf{$\tau_A$:} MLE power-law fit of the pooled energy-drop catalog with KS
test; compare to Gutenberg--Richter energy form.
\item \textbf{Transition order:} report the $\Pinf$ distribution and hysteresis vs.\
disorder; classify as continuous, first-order, or hybrid.
\item \textbf{Universality verdict:} assemble $\{\beta,\gamma,\nu,\tau,\tau_A\}$, check
the hyperscaling relation $2\beta+\gamma=d\nu$, and compare against 2D percolation,
random-fuse, fiber-bundle, and SOC reference values to decide whether VCEM defines a
distinct elastic crack-network class.
\end{itemize}

These outputs populate, in order, the six results subsections of the manuscript
(crack-network morphology; order parameters; susceptibilities and correlation length;
finite-size scaling and exponents; avalanche statistics; universality class).

%% ============================================================================
\section{Risks and contingencies}
\label{sec:risks}
%% ============================================================================

\begin{itemize}
\item \textbf{Large runs overrun the time budget.} Mitigation: the pyramid caps XL at
16 seeds; rule~4 spends the reserve on XXL \emph{only} if the exponents demand it.
\item \textbf{$\rhostar$ mis-set.} Mitigation: the pilot brackets the transition before
launch; the density sweep (C2) provides a fallback control parameter.
\item \textbf{Poor avalanche statistics.} Unlikely (thousands per run), but if a regime
is sparse the controller pools across seeds within that regime before fitting.
\item \textbf{Solver non-convergence on dense networks.} Flagged per run
(\texttt{converged\_flags}); non-converged runs are excluded from aggregates and
re-seeded, not silently averaged.
\item \textbf{Interruption / reboot.} The checkpointed agent loop resumes from
\texttt{manifest.csv} with no lost completed runs.
\end{itemize}

\end{document}
